% note: the official 2021 T3 solution PDF is a single page which ends after
% "Solution A" of part 3.2; the announced alternative method ("next") and
% part 3.3 are not present in the source file.
% note: the source (Bogota) uses decimal commas in some numbers; kept as
% printed.
\begin{description}
\item[(3.1)] Since the path between $A$ and $B$ is parabolic, the total
energy of the spacecraft is zero,
\[ E_{AB} = 0, \qquad \leftrightarrow \qquad \varepsilon = 1. \]
So, when you get to point B, then
\begin{equation*}
E = \tfrac{1}{2}mv_B^2 - \frac{GMm}{r_B} = 0; \tag*{[1 point]}
\end{equation*}
From where
\begin{equation*}
v_B = \sqrt{\frac{2GM}{r_B}} \tag*{[1 point]}
\end{equation*}
Replacing with available values
\begin{equation*}
v_B = \sqrt{\frac{2 \times (6{,}67 \times 10^{-11}) \times
  (6{,}4 \times 10^{23})}{6{,}8 \times 10^{6}}}
  \approx 3{,}55\ \mathrm{km\,s^{-1}} \tag*{[1 point]}
\end{equation*}

\item[(3.2)] There are two possible solutions

\textbf{Solution A}

Immediately after breaking % sic: "braking"
the total energy and the angular momentum change since the braking force is
tangential, but along the elliptical path BC the values with which it passes
through are conserved and take on during the entire journey. C; then applying
conservation of energy between points B$'$ (immediately after braking) and C
as well as conservation of the angular momentum between B$'$ and C, we have
\begin{gather*}
\tfrac{1}{2}mv_{B'}^2 - \frac{GMm}{r_B}
  = \tfrac{1}{2}mv_C^2 - \frac{GMm}{R_M} \quad Eq.~1
  \tag*{[1 point]} \\
L = mv_{B'}r_B = mv_C R_M \quad Eq.~2 \tag*{[1 point]}
\end{gather*}
From these two equations we solve for $v_C$ and it's result is used to
calculate the total energy in C which gives us:
\begin{equation*}
E_{elipse} = -\frac{GMm}{R_M + r_B} \approx -2.10 \times 10^{11}\,J
  \quad Eq.~3a \tag*{[2 points]}
\end{equation*}
In both methods, this and next, two points for getting the equation and two
points for calculation + negative sign + unit.
\end{description}
